\newif\ifglobal

%%%%%%%%%%%%%%%%%%%%%%%
%%% Compile to view %%%
%%%%%%%%%%%%%%%%%%%%%%%

%%%%%%%%%%%%%%%%%%%%%%%%%%%%%%%%%%%%%%%%%%%%%%%%%%%
%%% Readme:                                     %%%
%%% https://www.overleaf.com/read/bwdjvfjksvjy/ %%%
%%%%%%%%%%%%%%%%%%%%%%%%%%%%%%%%%%%%%%%%%%%%%%%%%%%

% >>> M?: marks a module setting number or important notice

% >>> M4: Set {./relative-path} to external files for this module <<<
\def \modulefiles {./17-SYSREG}
% To add an external file, use:
% (Replace '---' with actual filename)
% '\input{\modulefiles/tables/---.tex}' for tables
% '\includegraphics{\modulefiles/figures/---}' for figures
% '\subfile{\modulefiles/---}' for register files


% >>> M1: This module is in global project? <<<
% 'YES' - uncomment; 'NO' - comment out
\globaltrue

\ifglobal
    \documentclass[main\_\_EN.tex]{subfiles}
   
\else
    \documentclass[openany, 10pt]{book}

    % >>> M2: In ./00-shared/preamble-trm-module, also find
    %         settings for visibility of todo notes, watermark <<<
    \usepackage{preamble-trm-module}

    % >>> M3: Temporary labels in English,
    %         you can add your own labels there <<<
    \myexternaldocument{./00-shared/temp-labels-en}

\fi %\ifglobal


\setcounter{secnumdepth}{4}


% >>> M5: If this module needs any updates in
% './00-shared/preamble-trm-module.sty'
% (include additional package, etc.), contact your module PM
% or create the file './00-shared/changelog.tex' make the
% necessary updates yourself and document them in this file <<<

\raggedright


\begin{document}


% Sets language into which pre-defined bits of text are translated
\selectlanguage{english}

% Do not delete this block
\ifglobal
% Add the following front matter if
% this module is outside of global project
\else
    \InputIfFileExists{readme}{}{}
    {\let\clearpage\relax \listoftodos}
    {\let\clearpage\relax \listofdonetodos}
    \tableofcontents
    %\listoftables
    %\listoffigures
\fi


%%%%%%%%%%%%%%%%%%%%%%%%%
%%% TRM Module Begins %%%
%%%%%%%%%%%%%%%%%%%%%%%%%

\chapter{System Registers}
\label{mod:sysreg}
\hypertarget{sysreg}{}


\section{Overview}
\label{sec:sysreg-features}

\chipname{} supports the following auxiliary chip features:

\begin{itemize}
    \item External Memory Encryption/Decryption
    \item Anti-DPA attack security
    \item Bus timeout protection
\end{itemize}

Each auxiliary chip feature can be controlled with dedicated system registers. This chapter describes how to configure these system registers.


\section{Function Description}

\subsection{External Memory Encryption/Decryption Configuration}

\hyperref[regdesc:HPSYSTEMEXTERNALDEVICEENCRYPTDECRYPTCONTROLREG]{HP\_SYSTEM\_EXTERNAL\_DEVICE\_ENCRYPT\_DECRYPT\_CONTROL\_REG} configures encryption and decryption options of the external memory. For details about the External Memory Encryption and Decryption modules, please refer to Chapter \ref{mod:extmemencr}
\textit{\nameref{mod:extmemencr}}.


\subsection{Anti-DPA Attack Security Control}
\label{sec:sysreg-anti-dpa-attack-security-control}

\chipname{} has a dual protection mechanism against Differential Power Analysis (DPA) attacks at the hardware level. 

\begin{itemize}
    \item First, a mask mechanism is introduced in the symmetric encryption operation process, which interferes with the power consumption trajectory by masking the real data in the operation process. This security mechanism cannot be turned off.
    \item Second, the clock selected for the operation will change dynamically in real time, blurring the power consumption trajectory during the operation. For this security mechanism, \chipname{} provides three security levels for users to choose to adapt to different applications.
\end{itemize}


\begin{table}[H]
    \centering
    \caption{Security Level}
    \label{tab:sysreg-security-level}
    \begin{threeparttable}
    \begin{tabular}[c]{| l | c | c | c | c |}  
    \hline
    \rowcolor{lightgray}
        \textbf{Security-Level Name} & \textbf{Value} & \textbf{96M\_PLL\_CLK (MHz)} & \textbf{64M\_PLL\_CLK (MHz)} & \textbf{XTAL\_CLK (MHz)} \\ \hline
        SEC\_DPA\_OFF       & 0    & 96               & 64                 & 32                 \\\hline
        SEC\_DPA\_LOW       & 1 or 2    & (48,96]\tnote{A} & (32,64]\tnote{A}   & (16,32]\tnote{A}   \\\hline
        SEC\_DPA\_HIGH      & 3    & (32,96]\tnote{A} & (21.3,64]\tnote{A} & (10.6,32]\tnote{A}   \\\hline
    \end{tabular}
    \begin{tablenotes}
        \item[A] (x,y] means the operating frequency is greater than x MHz, and equal to or less than y MHz.
    \end{tablenotes}
    \end{threeparttable}
\end{table}

By default, the field \hyperref[fielddesc:HPSYSTEMSECDPALEVEL]{HP\_SYSTEM\_SEC\_DPA\_CFG\_SEL} in register \hyperref[regdesc:HPSYSTEMSECDPACONFREG]{HP\_SYSTEM\_SEC\_DPA\_CONF\_REG} is 0. In this case, the security-level is decided by the eFuse field \hyperref[fielddesc:EFUSESECDPALEVEL]{EFUSE\_SEC\_DPA\_LEVEL}. If the field \hyperref[fielddesc:HPSYSTEMSECDPALEVEL]{HP\_SYSTEM\_SEC\_DPA\_CFG\_SEL} is set to 1, the security-level is decided by \hyperref[fielddesc:HPSYSTEMSECDPALEVEL]{HP\_SYSTEM\_SEC\_DPA\_CFG\_LEVEL} in register \hyperref[regdesc:HPSYSTEMSECDPACONFREG]{HP\_SYSTEM\_SEC\_DPA\_CONF\_REG}.











\subsection{Bus Timeout Protection}

The Bus Timeout Protection function can be enabled and the timeout threshold can be configured through the configuration register. When a transfer is initiated, the counter inside the Timeout Protection module will increase by one every clock cycle. When the accumulated value is less than the timeout threshold and the bus receives a response from the slave, the internal counter is cleared. When the accumulated value is greater than the timeout threshold, if the slave device has not responded to the transfer, the Timeout Protection module will force the bus return signal to be pulled high. At the same time, it will report the interrupt and record the abnormal access address and master ID.


\subsubsection{CPU Peripheral Timeout Protection Register}

\hyperref[regdesc:HPSYSTEMCPUPERITIMEOUTCONFREG]{HP\_SYSTEM\_CPU\_PERI\_TIMEOUT\_CONF\_REG} is the timeout protection configuration register for accessing CPU peripheral registers.
CPU peripherals refer to the peripherals or modules whose addresses are in the range of 0x600C\_0000 \verb+~+ 0x600C\_FFFF. For corresponding peripheral information, please refer to Subsection \ref{para:sysmem-mod-peri-add-mapping} \textit{\nameref{para:sysmem-mod-peri-add-mapping}} in Chapter \ref{mod:sysmem} \textit{\nameref{mod:sysmem}}.

When a timeout occurs, the CPU\_PERI\_TIMEOUT\_INTR interrupt will be asserted.

\begin{itemize}
    \item \hyperref[regdesc:HPSYSTEMCPUPERITIMEOUTCONFREG]{HP\_SYSTEM\_CPU\_PERI\_TIMEOUT\_CONF\_REG}: Reserved.
    \item \hyperref[regdesc:HPSYSTEMCPUPERITIMEOUTADDRREG]{HP\_SYSTEM\_CPU\_PERI\_TIMEOUT\_ADDR\_REG}: 
    When a timeout occurs, this register will record the address of the timeout.
    \item \hyperref[regdesc:HPSYSTEMCPUPERITIMEOUTUIDREG]{HP\_SYSTEM\_CPU\_PERI\_TIMEOUT\_UID\_REG}: 
    When a timeout occurs, this register will record the Master-ID of the timeout.
\end{itemize}


\subsubsection{HP Peripheral Timeout Protection Register}

HP\_SYSTEM\_HP\_PERI\_TIMEOUT\_CONF\_REG is the timeout protection configuration register for accessing HP peripheral registers.
HP peripherals refer to the peripherals or modules whose addresses are in the range of 0x6000\_0000 \verb+~+ 0x6009\_FFFF. For corresponding peripheral information, please refer to Subsection \ref{para:sysmem-mod-peri-add-mapping} \textit{\nameref{para:sysmem-mod-peri-add-mapping}} in Chapter \ref{mod:sysmem} \textit{\nameref{mod:sysmem}}.

When a timeout occurs, the HP\_PERI\_TIMEOUT\_INTR interrupt will be asserted.

\begin{itemize}
    \item \hyperref[regdesc:HPSYSTEMHPPERITIMEOUTCONFREG]{HP\_SYSTEM\_HP\_PERI\_TIMEOUT\_CONF\_REG}: Reserved.
    \item \hyperref[regdesc:HPSYSTEMHPPERITIMEOUTADDRREG]{HP\_SYSTEM\_HP\_PERI\_TIMEOUT\_ADDR\_REG}: 
    When a timeout occurs, this register will record the address of the timeout.
    \item \hyperref[regdesc:HPSYSTEMHPPERITIMEOUTUIDREG]{HP\_SYSTEM\_HP\_PERI\_TIMEOUT\_UID\_REG}: 
    When a timeout occurs, this register will record the Master-ID of the timeout.
\end{itemize}

\subsubsection{LP Peripheral Timeout Protection Register}

\hyperref[regdesc:LPPERIBUSTIMEOUTCONFREG]{LP\_PERI\_BUS\_TIMEOUT\_CONF\_REG} is the timeout protection configuration register for accessing LP peripheral registers.
LP peripherals refer to the peripherals or modules whose addresses are in the range of 0x600B\_0000 \verb+~+ 0x600B\_FFFF. For corresponding peripheral information, please refer to Subsection \ref{para:sysmem-mod-peri-add-mapping} \textit{\nameref{para:sysmem-mod-peri-add-mapping}} in Chapter \ref{mod:sysmem} \textit{\nameref{mod:sysmem}}.

When a timeout occurs, the LP\_PERI\_TIMEOUT\_INTR interrupt will be asserted.

\begin{itemize}
    \item \hyperref[regdesc:LPPERIBUSTIMEOUTCONFREG]{LP\_PERI\_BUS\_TIMEOUT\_CONF\_REG}: Reserved.
    \item \hyperref[regdesc:LPPERIBUSTIMEOUTADDRREG]{LP\_PERI\_BUS\_TIMEOUT\_ADDR\_REG}: 
    When a timeout occurs, this register will record the address of the timeout.
    \item \hyperref[regdesc:LPPERIBUSTIMEOUTUIDREG]{LP\_PERI\_BUS\_TIMEOUT\_UID\_REG}: 
    When a timeout occurs, this register will record the Master-ID of the timeout.
\end{itemize}





%%%%%%%%%%%%%%%%%%%%%%%%
%%% Register Summary %%%
%%%%%%%%%%%%%%%%%%%%%%%%

% >>> Update the sample text below and insert
%     the info on register summary into the subfile <<<

\clearpage
\begin{landscape}
\hypertarget{sysreg-reg-summ}{}
\section{Register Summary}
\label{sec:sysreg-reg-summ}

The addresses related of LP Peripheral timeout registers are relative to Low-power Peripheral base address provided in Table \ref{tab:sysmem-base-address} in Chapter \ref{mod:sysmem} \textit{\nameref{mod:sysmem}}, and the others addresses in this section are relative to System Registers base address provided in Table \ref{tab:sysmem-base-address} in Chapter \ref{mod:sysmem} \textit{\nameref{mod:sysmem}}.

The abbreviations given in Column \textbf{Access} are explained in Section \hyperref[glossary-access-types]{\textit{Access Types for Registers}}.

\subfile{\modulefiles/17-SYSREG-reg-summ__EN}
\end{landscape}

%%%%%%%%%%%%%%%%%
%%% Registers %%%
%%%%%%%%%%%%%%%%%

% >>> Update the sample text below and insert
%      the info on registers into the subfile <<<

\clearpage
\section{Registers}

The addresses related of LP Peripheral timeout registers are relative to Low-power Peripheral base address provided in Table \ref{tab:sysmem-base-address} in Chapter \ref{mod:sysmem} \textit{\nameref{mod:sysmem}}, and the others addresses in this section are relative to System Registers base address provided in Table \ref{tab:sysmem-base-address} in Chapter \ref{mod:sysmem} \textit{\nameref{mod:sysmem}}.

\subfile{\modulefiles/17-SYSREG-reg__EN}

\end{document}
